% Options for packages loaded elsewhere
\PassOptionsToPackage{unicode}{hyperref}
\PassOptionsToPackage{hyphens}{url}
\documentclass[
  brazilian,
  12pt,
  a4paper,
]{article}
\usepackage{xcolor}
\usepackage{amsmath,amssymb}
\setcounter{secnumdepth}{-\maxdimen} % remove section numbering
\usepackage{iftex}
\ifPDFTeX
  \usepackage[T1]{fontenc}
  \usepackage[utf8]{inputenc}
  \usepackage{textcomp} % provide euro and other symbols
\else % if luatex or xetex
  \usepackage{unicode-math} % this also loads fontspec
  \defaultfontfeatures{Scale=MatchLowercase}
  \defaultfontfeatures[\rmfamily]{Ligatures=TeX,Scale=1}
\fi
\usepackage{lmodern}
\ifPDFTeX\else
  % xetex/luatex font selection
  \setmainfont[]{TeX Gyre Termes}
  \setsansfont[]{TeX Gyre Heros}
  \setmonofont[]{DejaVu Sans Mono}
\fi
% Use upquote if available, for straight quotes in verbatim environments
\IfFileExists{upquote.sty}{\usepackage{upquote}}{}
\IfFileExists{microtype.sty}{% use microtype if available
  \usepackage[]{microtype}
  \UseMicrotypeSet[protrusion]{basicmath} % disable protrusion for tt fonts
}{}
\usepackage{setspace}
\makeatletter
\@ifundefined{KOMAClassName}{% if non-KOMA class
  \IfFileExists{parskip.sty}{%
    \usepackage{parskip}
  }{% else
    \setlength{\parindent}{0pt}
    \setlength{\parskip}{6pt plus 2pt minus 1pt}}
}{% if KOMA class
  \KOMAoptions{parskip=half}}
\makeatother
\usepackage{longtable,booktabs,array}
\usepackage{calc} % for calculating minipage widths
% Correct order of tables after \paragraph or \subparagraph
\usepackage{etoolbox}
\makeatletter
\patchcmd\longtable{\par}{\if@noskipsec\mbox{}\fi\par}{}{}
\makeatother
% Allow footnotes in longtable head/foot
\IfFileExists{footnotehyper.sty}{\usepackage{footnotehyper}}{\usepackage{footnote}}
\makesavenoteenv{longtable}
\ifLuaTeX
\usepackage[bidi=basic,shorthands=off,]{babel}
\else
\usepackage[bidi=default,shorthands=off,]{babel}
\fi
\ifPDFTeX
\else
\babelfont{rm}[]{TeX Gyre Termes}
\fi
\ifLuaTeX
  \usepackage{selnolig} % disable illegal ligatures
\fi
\setlength{\emergencystretch}{3em} % prevent overfull lines
\providecommand{\tightlist}{%
  \setlength{\itemsep}{0pt}\setlength{\parskip}{0pt}}
\usepackage{graphicx}
\usepackage{array}
\usepackage{etoolbox}
\AtBeginEnvironment{tabular}{\footnotesize}
\AtBeginEnvironment{longtable}{\footnotesize}
\usepackage{bookmark}
\IfFileExists{xurl.sty}{\usepackage{xurl}}{} % add URL line breaks if available
\urlstyle{same}
\hypersetup{
  pdflang={pt-BR},
  pdfkeywords={Armadilha da renda média, IDEB, Sertão de
Crateús, Ceará, Economia da educação, Painel municipal},
  hidelinks,
  pdfcreator={LaTeX via pandoc}}

\author{}
\date{}


% ---------- Formatação ABNT / RBEP ----------
\usepackage[a4paper,left=3cm,right=2cm,top=3cm,bottom=2cm]{geometry}
\usepackage{indentfirst}
\setlength{\parindent}{1.25cm}
\setlength{\parskip}{4pt plus 1pt}
\usepackage{titlesec}
\titleformat{\section}{\normalfont\bfseries}{\thesection}{0.6em}{}
\titleformat{\subsection}{\normalfont\bfseries}{\thesubsection}{0.6em}{}
\titleformat{\subsubsection}{\normalfont\bfseries\itshape}{\thesubsubsection}{0.6em}{}
\titlespacing*{\section}{0pt}{18pt plus 4pt minus 2pt}{8pt plus 2pt}
\titlespacing*{\subsection}{0pt}{14pt plus 3pt minus 2pt}{6pt plus 1pt}
\usepackage[font=small,labelsep=endash,justification=centering]{caption}
\captionsetup[figure]{position=below}
\usepackage{fancyhdr}
\pagestyle{fancy}
\fancyhf{}
\fancyhead[R]{\thepage}
\renewcommand{\headrulewidth}{0pt}
\fancypagestyle{plain}{\fancyhf{}\fancyhead[R]{\thepage}\renewcommand{\headrulewidth}{0pt}}
\usepackage{etoolbox}
\AtBeginEnvironment{longtable}{\footnotesize}
\AtBeginEnvironment{tabular}{\footnotesize}
\setlength{\tabcolsep}{4pt}
\renewcommand{\arraystretch}{1.15}

\begin{document}

\setstretch{1.5}
\begin{center}
{\Large\bfseries Educação, Armadilha da Renda Média e Desempenho Escolar:}\\[4pt]
{\large\bfseries Evidências do IDEB no Sertão de Crateús e no Ceará (2005–2025)\\ como Ruptura da Estagnação}
\end{center}

\begin{center}
\small Prof. Marcelo Claro --- Mestre em Educação, Pesquisador em Economia da Educação\\
\small ORCID: \href{https://orcid.org/0000-0001-8996-2887}{0000-0001-8996-2887}
\end{center}

\begin{quote}
\textbf{Fontes e proveniência:} todos os dados são de fontes oficiais
auditáveis (INEP; IBGE; World Bank; OCDE; FAO; USPTO; CONAB; IPEA). Cada
download está registrado em manifesto próprio com URL, data-hora e hash
SHA-256. Scripts, saídas completas e o relatório independente de
validação por banca simulada acompanham o material suplementar (Apêndice
A).
\end{quote}

\subsection{Resumo}\label{resumo}

Este artigo articula duas escalas para examinar se a educação constitui
fator de ruptura da armadilha da renda média no Brasil. Na escala macro
(1960--2026; 11 países; World Bank, IBGE, OCDE-PISA, FAO, USPTO), seis
de sete correlações de Pearson são significativas: educação--PIB
(r=0,86), PISA--PIB (r=0,95), internet--PIB (r=0,66--0,71) e
Gini--educação (r=−0,94 a −0,95); o PIB--insegurança alimentar não é
significativo (r=−0,39; p=0,155). O Brasil acumula 14 anos abaixo do
pico de PIB per capita (US\$ 13.397, em 2011). Na escala micro (Sertão
de Crateús, CE; nove municípios; painel 2013--2023, n=108; inferência
com bootstrap por cluster G=9, erros clusterizados com correção de
pequena amostra e \emph{wild cluster bootstrap} Rademacher de 9.999
réplicas), não se detecta associação entre renda municipal e IDEB nem
entre municípios (r=−0,24; IC95\% {[}−0,77; 0,50{]}; p=0,52) nem dentro
deles (β=2,61; p=0,20; p \emph{wild}=0,23), com poder insuficiente
(MDES=6,01 por unidade de log; 0,57 ponto por +10\% de PIB) e
equivalência não estabelecida (TOST ±0,5: p=0,85). Registram-se 107/108
metas do INEP cumpridas no mesmo ano (99,1\%) e ganho médio de 5,93
pontos (2007--2025) sem associação com renda (r=−0,14; p=0,73).
Conclui-se que a política educacional coordenada (PAIC/CE) desacoplou
localmente o avanço escolar da renda municipal --- ruptura micro
compatível com estagnação macro --- sem reivindicação causal.

\textbf{Palavras-chave:} armadilha da renda média; IDEB; Sertão de
Crateús; Ceará; economia da educação; painel municipal; \emph{wild
cluster bootstrap}; regime de colaboração.

\subsection{Abstract}\label{abstract}

This article links two scales to examine whether education breaks
Brazil's middle-income trap. At the macro scale (1960--2026; 11
countries), six of seven Pearson correlations are significant:
education--GDP (r=0.86), PISA--GDP (r=0.95), internet--GDP
(r=0.66--0.71), Gini--education (r=−0.94 to −0.95); GDP--food insecurity
is not (r=−0.39; p=0.155). Brazil has spent 14 years below its 2011
GDP-per-capita peak (US\$13,397). At the micro scale (Sertão de Crateús,
CE; nine municipalities; panel 2013--2023, n=108; cluster-bootstrap
inference with G=9, small-sample-corrected CRVE and Rademacher wild
cluster bootstrap with 9,999 replications), no income--IDEB association
is detected across municipalities (r=−0.24; 95\% CI {[}−0.77; 0.50{]};
p=0.52) nor within them (β=2.61; p=0.20; wild p=0.23), with insufficient
power (MDES=6.01 per log unit; 0.57 point per +10\% GDP) and equivalence
not established (TOST ±0.5: p=0.85). We record 107/108 INEP targets met
in the same year (99.1\%) and a mean gain of 5.93 points (2007--2025)
unrelated to income (r=−0.14; p=0.73). Coordinated policy (PAIC/CE)
locally decoupled schooling gains from municipal income --- a micro
rupture consistent with macro stagnation --- without causal claims.

\textbf{Keywords:} middle-income trap; IDEB; Sertão de Crateús; Ceará;
economics of education; municipal panel; wild cluster bootstrap;
collaboration regime.

\subsection{1. Introdução: o enigma brasileiro em duas
escalas}\label{introduuxe7uxe3o-o-enigma-brasileiro-em-duas-escalas}

Em 2011, o PIB per capita do Brasil atingiu US\$ 13.397 --- o maior
valor da série histórica iniciada em 1960. Nos catorze anos seguintes, a
economia brasileira não conseguiu recuperar esse patamar: a média de
2012--2024 situou-se em US\$ 9.855, isto é, 26,4\% abaixo do pico, com
trajetória marcada pela contração de cerca de um terço entre 2011 e
2016, recuperação parcial até 2020, novo golpe pandêmico e retomada
incompleta. Esse platô prolongado é a manifestação mais visível do
fenômeno que a literatura internacional denominou \textbf{armadilha da
renda média}: a interrupção sistemática do crescimento antes que a
economia complete a transição de um padrão intensivo em mão de obra
barata para outro sustentado por inovação, produtividade e serviços
complexos (Gill; Kharas, 2007; Aiyar et al., 2013; Cherif; Hasanov,
2015).

O caso brasileiro reúne, com rara clareza, os sintomas descritos nessa
literatura. Primeiro, a matriz exportadora voltou-se crescentemente às
commodities, com baixo conteúdo tecnológico doméstico. Segundo, a
produtividade agregada permaneceu praticamente estagnada ao longo do
período. Terceiro, os indicadores de esforço inovativo --- depósitos de
patentes e densidade de atividades intensivas em pesquisa --- oscilaram
sem tendência consistente de convergência com as economias de fronteira.
O contraste é ainda mais intrigante porque coincidiu com avanços sociais
e digitais relevantes: a adoção de inteligência artificial pelas
indústrias brasileiras saltou de 16,9\% (2022) para 41,9\% (2024),
segundo a PINTEC/IBGE; o país saiu do Mapa da Fome, com a população em
insegurança alimentar grave recuando de 70,3 para 28,5 milhões; e o
coeficiente de Gini atingiu 50,3 em 2024, o menor nível desde 1980. Se
progresso social, conectividade e adoção tecnológica avançaram
simultaneamente, por que o salto de renda não ocorreu?

A tese aqui examinada é que \textbf{a qualidade da educação é a variável
estrutural mais fortemente associada ao desenvolvimento econômico
brasileiro, mas é condição necessária e não suficiente}. Na dimensão
temporal longa (1960--2026), a correlação entre anos de escolaridade e
PIB per capita alcança r=0,86, e a correlação entre desempenho no PISA
--- medida de proficiência, e não de frequência --- e PIB per capita
chega a r=0,95 na comparação internacional. A desigualdade, por sua vez,
associa-se negativa e fortemente à educação (Gini--educação r=−0,94 a
−0,95), sugerindo que a distribuição das oportunidades escolares é parte
integrante do mecanismo de desenvolvimento. Nenhuma dessas variáveis,
tomada isoladamente, rompe a armadilha: elas formam um ecossistema no
qual a educação opera como condição habilitante, mas cuja efetivação
depende de complementaridades --- conectividade, capacidade produtiva de
absorver habilidades, segurança alimentar e menor dispersão de renda.

Resta, contudo, uma pergunta decisiva que a literatura macro raramente
enfrenta: \textbf{se a educação é o fator de ruptura, municípios pobres
podem avançar educacionalmente sem esperar pelo ciclo econômico local?}
A pergunta importa porque, sob a leitura determinista, gestores de
municípios de baixa renda deveriam aguardar crescimento econômico para
só então colher resultados escolares --- posição com consequências
práticas profundas em regiões de pobreza estrutural, como o Semiárido
nordestino. A alternativa institucionalista sustenta que políticas
pedagógicas coordenadas podem produzir aprendizagem independentemente do
ciclo econômico municipal, desde que existam arranjos estáveis de
colaboração federativa, formação docente continuada, materiais
estruturados e regimes regulares de avaliação.

Para submeter essas leituras concorrentes a um teste empírico
delimitado, este artigo concentra-se no \textbf{Sertão de Crateús}
(microrregião IBGE 23018, Ceará): nove municípios pequenos, de baixa
renda per capita, situados no Semiárido, sob a mesma coordenação
estadual do Programa Alfabetização na Idade Certa (PAIC), com série
completa do Índice de Desenvolvimento da Educação Básica (IDEB) de 2005
a 2025 na rede municipal e indicadores econômicos oficiais (PIB
municipal deflacionado pelo IPCA e rendimentos do Censo Demográfico
2022). O contexto é analiticamente privilegiado porque combina
homogeneidade institucional (mesma política estadual), heterogeneidade
econômica moderada mas informativa e disponibilidade de metas públicas
bianuais pactuadas pelo INEP --- o que permite avaliar não apenas níveis
e tendências, mas também o cumprimento de compromissos explícitos de
desempenho.

\textbf{Pergunta de pesquisa (RQ).} O artigo responde a três questões
encadeadas: (i) na escala macro, quais associações empíricas sustentam
--- ou qualificam --- a educação como fator de ruptura da armadilha da
renda média?; (ii) na microrregião de Crateús, a renda municipal
associa-se ao IDEB \textbf{entre municípios} (transversalmente) e
\textbf{dentro de cada município} (temporalmente)?; (iii) o ganho
educacional observado depende da renda municipal, e as metas pactuadas
junto ao INEP foram cumpridas nos anos de referência?

\textbf{Contribuições.} O estudo oferece três contribuições distintas.
Primeira, integra num mesmo desenho as duas escalas usualmente tratadas
separadamente --- a macroeconomia do desenvolvimento (armadilha da renda
média) e a microeconomia da política educacional municipal --- com
triangulação explícita entre microrregião, estado e país. Segunda,
aplica à escala municipal um protocolo inferencial conservador raramente
adotado nesse nível de agregação: bootstrap por clusters municipais,
erros-padrão clusterizados com correção de pequena amostra, \emph{wild
cluster bootstrap} sob a hipótese nula, validação interna por exclusão
de municípios (LOOCV), dimensionamento explícito de poder (MDES) e
testes de equivalência (TOST) com margem substantiva pré-especificada.
Terceira, documenta com proveniência criptográfica (hashes SHA-256 de
cada arquivo-fonte) todo o acervo de dados, permitindo auditoria
independente ponta a ponta --- requisito metodológico ainda pouco
difundido em pesquisas educacionais municipais.

Deliberadamente, o desenho é \textbf{observacional e associativo}: não
há variação exógena de renda que autorize linguagem causal. O termo
``efeito'' aparece apenas como sinônimo técnico de coeficiente de
regressão condicional. Os limites inferenciais são declarados a priori e
operacionalizados por critérios de falsificabilidade (Seção 3.6), e toda
a cadeia de produção --- coleta, limpeza, estimação e saídas --- é
reprodutível com sementes aleatórias fixas.

O restante do artigo organiza-se assim: a Seção 2 revisita a
fundamentação teórica e empírica sobre armadilha da renda média, o IDEB
como instrumento de política e o caso cearense; a Seção 3 detalha dados,
variáveis, hipóteses e procedimentos inferenciais; a Seção 4 apresenta
os resultados nas duas escalas, com tabelas e figuras; a Seção 5 discute
interpretações, mecanismos e implicações de política; a Seção 6
consolida limitações; a Seção 7 conclui e aponta agenda futura. O
objetivo geral é avaliar se a renda municipal condiciona o desempenho
escolar; a metodologia foi desenhada para tornar essa pergunta
falsificável em cada escala analisada.

\subsection{2. Fundamentação}\label{fundamentauxe7uxe3o}

\subsubsection{2.1 Armadilha da renda média e educação como
ruptura}\label{armadilha-da-renda-muxe9dia-e-educauxe7uxe3o-como-ruptura}

A noção de armadilha da renda média foi sistematizada por Gill e Kharas
(2007) ao observarem que economias do Leste Asiático que haviam crescido
rapidamente com base em manufatura intensiva em trabalho enfrentavam
dificuldades crescentes para sustentar o ritmo quando os salários
convergiam para patamares intermediários. A formulação ganhou tratamento
econométrico com Aiyar et al.~(2013), que identificaram probabilidades
elevadas de desaceleração em faixas específicas de renda per capita, e
prescrições de política com Cherif e Hasanov (2015), para quem a escape
exige ``políticas de verdade'' de industrialização complexa e inovação,
e não atalhos de curto prazo. Três regularidades emergem dessa tradição:
(i) a transição exige capacidades de absorção --- força de trabalho
capaz de aprender, adaptar e melhorar processos; (ii) instituições que
coordenem investimento privado e formação de habilidades reduzem a
probabilidade de travamento; (iii) a qualidade --- e não apenas a
quantidade --- de escolaridade é o elo entre sistema educacional e
crescimento.

Essa última regularidade foi formalizada por Hanushek e Woessmann
(2015): o que prediz crescimento de longo prazo entre países é o
\textbf{capital cognitivo} --- proficiência medida em avaliações
internacionais padronizadas --- e não os anos de assento escolar. Para o
Brasil, a distinção é crucial. O país quase universalizou o acesso e
ampliou fortemente os anos médios de estudo, mas permanece distante das
médias da OCDE em proficiência (abaixo de 400 pontos em matemática no
PISA, contra média de 472). A correlação empírica documentada neste
artigo --- PISA--PIB r=0,95 na amostra de 11 países (1960--2026) ---
reforça que a variável relevante para romper a armadilha é aprendizagem
efetiva, o que reposiciona o IDEB, indicador que combina fluxo e
proficiência, como objeto central de análise para o caso nacional.

Por fim, a literatura recente sobre a armadilha enfatiza que nenhum
fator isolado é suficiente. A experiência comparada sugere ecossistemas:
educação básica de qualidade, conectividade digital, ambiente inovativo,
segurança alimentar e menor concentração de renda atuam conjuntamente. É
exatamente essa hipótese de \textbf{ecossistema} que a parte macro deste
artigo operationaliza por meio de sete correlações de longo prazo, lidas
como descrição de regularidades --- nunca como identificação causal.

\subsubsection{2.2 IDEB, metas e accountability: promessas e
riscos}\label{ideb-metas-e-accountability-promessas-e-riscos}

Criado em 2007, o IDEB sintetiza, em um único número por escola,
município, estado e rede, duas dimensões: fluxo escolar (taxa de
aprovação) e proficiência média (SAEB), com metas bianuais projetadas
pelo INEP para horizonte de duas décadas (Soares; Xavier, 2013). Essa
arquitetura transformou o indicador em peça central de um regime de
responsabilização educacional: metas públicas, comparáveis e datadas
criam referencial de cobrança social e gerencial. A literatura
avaliativa reconhece sua função mobilizadora --- dar visibilidade ao
desempenho e induzir ação coordenada das redes (Lacruz; Américo;
Carniel, 2019) ---, mas alerta para riscos clássicos de sistemas de alta
consequência: estreitamento curricular em direção ao que é testado,
manipulação de fluxo e ensinamento para o teste (Afonso, 2012;
Schneider; Nardi, 2014). Evidência internacional indica que avaliações
com consequências associam-se a ganhos mensuráveis em matemática
(Brunello; Kiss, 2022), o que torna plausível que a combinação
brasileira de medida pública e meta explícita altere comportamentos
escolares --- questão que permanece aberta quanto à sua sustentabilidade
e distribuição de ganhos.

Do ponto de vista deste artigo, o IDEB tem duas virtudes metodológicas
decisivas para o teste municipal proposto. Primeiro, provém uma
\textbf{métrica comum e longitudinal} a todas as redes, permitindo
painéis balanceados por etapa e unidade federada. Segundo, as
\textbf{metas projetadas pelo INEP} oferecem um critério externo de
desempenho --- observado versus projetado no mesmo ano --- que independe
da renda local, viabilizando a hipótese H3 (cumprimento universal de
metas em região pobre) sem circularidade analítica.

\subsubsection{2.3 Renda e desempenho: associação transversal versus
temporal}\label{renda-e-desempenho-associauxe7uxe3o-transversal-versus-temporal}

A relação renda--desempenho escolar está entre as mais replicadas da
pesquisa educacional, mas sua interpretação exige uma distinção
analítica frequentemente negligenciada. Na dimensão \textbf{transversal}
--- entre unidades em um mesmo momento --- municípios mais ricos tendem,
em média, a apresentar IDEB mais alto, reflexo de dotações
socioeconômicas das famílias e da capacidade fiscal das redes (Duarte,
2013; Andrews; Vries, 2012). Na dimensão \textbf{temporal} --- dentro de
cada unidade ao longo do tempo ---, a covariação entre mudanças de renda
e mudanças de desempenho é substancialmente mais fraca e instável: redes
inteiras avançam em períodos nos quais a renda local mal se move,
sugerindo que fatores pedagógico-institucionais dominam a variância
intra-unidade (Andrews; Vries, 2012). Confundir as duas dimensões
alimenta a falácia ecológica: atribuir à renda local um papel
explicativo que só existe na comparação agregada entre unidades muito
distintas.

Metodologicamente, quando a unidade de clusterização é o município e o
número de clusters é pequeno (neste estudo, G=9), a inferência
convencional superestima precisão. As recomendações consolidadas incluem
bootstrap por cluster com reamostragem de unidades inteiras (Cameron;
Gelbach; Miller, 2008), correções de pequena amostra nos erros-padrão
clusterizados, distribuição t com G−1 graus de liberdade e, como
procedimento decisivo, o \emph{wild cluster bootstrap} impondo a
restrição nula (MacKinnon; Webb, 2017). Este artigo adota integralmente
esse protocolo e acrescenta dois elementos ainda menos comuns na
literatura educacional municipal: dimensionamento de poder (MDES) e
testes de equivalência (TOST) com margem substantiva, impedindo que
ausência de significância seja convertida retoricamente em prova de
nulidade.

\subsubsection{2.4 O Semiárido cearense como caso
crítico}\label{o-semiuxe1rido-cearense-como-caso-cruxedtico}

Regiões semiáridas combinam pobreza estrutural, adversidade climática e
historicamente baixos indicadores sociais, o que as tornava candidatas
canônicas ao determinismo socioeconômico: onde a renda é baixa e
volátil, esperar-se-ia desempenho escolar igualmente deprimido. O Ceará
desafiou essa expectativa a partir de 2007 com o PAIC (Lei Estadual nº
14.026/2007), regime formal de colaboração entre estado e municípios
centrado na alfabetização na idade certa: rede estadual de formação de
professores, materiais didáticos estruturados, assessment anual externo
de fluência leitora no 2º ano e governança compartilhada por convênios
de adesão. Avaliações independentes documentaram ganhos expressivos de
alfabetização e aprendizagem (Costa; Carnoy, 2015; Cruz; Ribeiro;
Batista, 2022), e análises intranacionais destacaram o Ceará entre os
casos de maior progresso concentrado nos anos iniciais das redes
municipais (Carnoy et al., 2017; Segatto; Oliveira; Silva, 2024).

O Sertão de Crateús, nesse arranjo, funciona como teste de estresse: se
mesmo em municípios de renda entre as mais baixas do estado o IDEB
converge para patamares altos e cumpre metas, então a hipótese do
desacoplamento entre política pedagógica coordenada e ciclo econômico
local ganha suporte empírico direto na escala em que a política de fato
acontece --- o município.

\subsection{3. Método}\label{muxe9todo}

\subsubsection{3.1 Delimitação territorial e fontes de
dados}\label{delimitauxe7uxe3o-territorial-e-fontes-de-dados}

A unidade territorial primária é a microrregião geográfica IBGE 23018
--- Sertão de Crateús --- composta por nove municípios: Ararendá
(2301257), Crateús (2304103), Independência (2305605), Ipaporanga
(2305654), Monsenhor Tabosa (2308609), Nova Russas (2309300), Novo
Oriente (2309409), Quiterianópolis (2311264) e Tamboril (2313203). A
escolha obedece a três critérios declarados \emph{a priori}:
homogeneidade institucional (todos os municípios aderiram ao PAIC desde
os primeiros ciclos), completude da série IDEB na rede municipal e
disponibilidade de renda oficial por município em todo o período de
interesse.

As fontes primárias, todas oficiais e publicamente auditáveis,
compreendem: (i) \textbf{INEP/IDEB}, edição de 2025, série 2005--2025,
redes Municipal e Pública, anos iniciais e finais do ensino fundamental,
com metas projetadas bienais; (ii) \textbf{IBGE --- PIB dos Municípios},
anos 2010--2021, PIB per capita a preços correntes; (iii) \textbf{IBGE
--- Censo Demográfico 2022}, agregados por município de rendimento
nominal mensal médio (V06004) e mediano (V06006) das pessoas
responsáveis por domicílios, além de totais de moradores; (iv)
\textbf{IPEA Data}, série PRECOS12/IPCA12 (número-índice do IPCA, base
dezembro de 1993=100; fonte original IBGE/SNIPC), com 1.680 registros
mensais, utilizada na construção do deflator; (v) \textbf{IBGE --- Malha
Municipal 2023} para cartografia; e, para a análise macro, (vi)
\textbf{World Bank Open Data}, \textbf{OCDE/PISA}, \textbf{FAO},
\textbf{USPTO} e \textbf{CONAB}. Cada download foi registrado em
manifesto de proveniência com URL, data-hora e hash SHA-256
(\texttt{data/raw/SOURCE\_MANIFEST.json}), garantindo rastreabilidade
criptográfica.

Um ajuste de escopo merece nota: a demanda original mencionava o índice
comercial IPC Marketing (IPCE) como proxy local de renda; por não
possuir acesso público auditável, ele foi substituído --- com a
limitação declarada --- por duas medidas oficiais complementares: PIB
per capita real (dimensão produtiva, painel anual) e rendimento médio do
responsável (dimensão domiciliar, corte censitário 2022). A substituição
preserva a finalidade analítica e eleva a auditabilidade.

\subsubsection{3.2 Variáveis e construção do
painel}\label{variuxe1veis-e-construuxe7uxe3o-do-painel}

\textbf{Variável de resultado.} Nota do IDEB municipal por etapa (anos
iniciais --- AI; anos finais --- AF), rede pública Municipal, edições
bienais de 2005 a 2025. Como o IDEB combina aprovação e proficiência
padronizada, variações refletem tanto fluxo quanto aprendizagem; a
leitura conjunta das duas componentes não é explorada aqui por
indisponibilidade de decomposição municipal auditável em todos os anos.

\textbf{Variável explicativa.} Logaritmo natural do PIB per capita
municipal \textbf{deflacionado} para reais de 2021. O deflator foi
construído a partir do IPCA médio anual: para cada ano \emph{a},
calculou-se a média aritmética dos índices mensais e o fator de
conversão \texttt{fator\_a~=~IPCA\_2021~/~IPCA\_a}; o valor nominal foi
multiplicado pelo fator do respectivo ano. O logaritmo confere
interpretação de elasticidade --- um acréscimo de 10\% no PIB per capita
real corresponde a Δln ≈ 0,0953 --- e simetria com a literature empírica
de crescimento.

\textbf{Estratégia de emparelhamento temporal.} Cada observação do IDEB
no ano \emph{t} foi pareada ao log do PIB per capita real de \emph{t−2}.
A defasagem de dois anos, declarada \emph{a priori} como principal,
aproxima a janela de exposição econômica da coorte avaliada (o PIB do
próprio ano da prova não pode influenciar desempenho já realizado) e
mitiga contemporaneidade mecânica. O painel resultante para a
microrregião cobre os anos de IDEB 2013--2023 (seis ondas bianuais), com
estrutura 9 municípios × 6 anos × 2 etapas = 108 células-alvo e apenas
uma ausência de IDEB (missingness de 0,5\%; n=108 após listwise).
Repetiram-se todos os procedimentos para o Ceará (184 municípios;
missing 1,3\%; n=2.202) e para o Brasil (5.451 municípios na base;
missing 26,2\%, decorrente majoritariamente de redes não municipais e
lacunas de série; n=48.695), como triangulação de escalas.

\textbf{Medidas de contexto.} Rendimento médio mensal das pessoas
responsáveis por domicílio (Censo 2022) entrou exclusivamente nas
análises do ganho educacional (H3b) e na caracterização socioeconômica
municipal, por ser a única fonte censitária domiciliar disponível
simultaneamente para os nove municípios no período recente.

\subsubsection{3.3 Hipóteses}\label{hipuxf3teses}

As hipóteses foram registradas antes da execução analítica e mapeiam-se
biunivocamente aos estimadores da Seção 3.4:

\begin{itemize}
\item
  \textbf{H1 (associação transversal).} A correlação de Pearson entre
  médias temporais municipais de IDEB e de log do PIB per capita real é
  positiva e discernível de zero (r ≥ 0,3; IC95\% exclui zero), estimada
  sobre G=9 médias municipais com incerteza por bootstrap de clusters.
\item
  \textbf{H2 (associação temporal fraca ou nula).} Dentro das unidades,
  primeiras diferenças apresentam correlações próximas de zero e o
  coeficiente de efeitos fixos de município×etapa e ano é
  estatisticamente indistinguível de zero sob erros clusterizados com
  correção de pequena amostra e sob \emph{wild cluster bootstrap}.
  Trata-se de hipótese avaliada por \textbf{não-detecção qualificada}:
  somente com MDES e TOST reportados é possível classificar
  adequadamente um resultado não significativo.
\item
  \textbf{H3a (metas).} Em cada par (município, etapa, ano t nos anos
  \{2011, 2013, 2015, 2017, 2019, 2021\}), o IDEB observado é maior ou
  igual à meta projetada pelo INEP para aquele mesmo ano --- critério de
  cumprimento sincronizado, mais exigente que comparações defasadas.
\item
  \textbf{H3b (ganho e renda).} O ganho de IDEB em AI entre 2007 e 2025
  não se associa à renda municipal censitária (correlação compatível com
  zero, IC95\% amplo reportado).
\end{itemize}

\subsubsection{3.4 Estratégia de inferência (protocolo
conservador)}\label{estratuxe9gia-de-inferuxeancia-protocolo-conservador}

Com apenas G=9 clusters, cada procedimento foi escolhido para minimizar
super-rejeição e maximizar transparência. A sequência completa é:

\begin{enumerate}
\def\labelenumi{\arabic{enumi}.}
\item
  \textbf{Between-estimator (H1).} Para cada município calculou-se a
  média do IDEB e do log do PIB real sobre as ondas do painel;
  reporta-se a correlação de Pearson entre essas médias, com intervalo
  de 95\% e valor-p obtidos por \textbf{bootstrap de clusters} ---
  sorteio de G municípios com reposição, 5.000 réplicas, semente fixa
  42, percentis 2,5--97,5 e proporção de réplicas com inversão de sinal
  como p aproximado. Jackknife por município acompanha como diagnóstico
  de influência.
\item
  \textbf{Correlações pooled.} Pearson sobre todas as observações do
  painel, com IC e p por bootstrap de clusters (5.000 réplicas na micro;
  500 nas escalas com G\textgreater50), reportados como complemento
  descritivo --- o tamanho efetivo transversal permanece sendo o número
  de municípios.
\item
  \textbf{Primeiras diferenças.} Correlação de Pearson sobre diferenças
  de primeira ordem dentro de cada município×etapa, eliminando
  heterogeneidade permanente; p nominal reportado com advertência de
  poder reduzido.
\item
  \textbf{Efeitos fixos (H2).} Regressão de IDEB em ln PIB real t−2 com
  efeitos fixos duplos de \textbf{município×etapa} (18 grupos) e de
  \textbf{ano} (5 dummies), estimada por within-transformation.
  Reportam-se: erro-padrão homocedástico (transparência),
  \textbf{erro-padrão clusterizado por município com correção de pequena
  amostra} (fator G/(G−1)·(N−1)/(N−K)) e teste t com G−1=8 graus de
  liberdade, além do \textbf{wild cluster bootstrap de Rademacher sob H0
  com 9.999 réplicas}, que impõe a nulidade ao reamostrar sinais dos
  resíduos restritos por cluster --- procedimento recomendado
  precisamente para G\textless40 (MacKinnon; Webb, 2017).
\item
  \textbf{Validação interna por exclusão (LOOCV).} Para cada município
  \emph{m}, recalculou-se a correlação de níveis usando as demais
  observações e reportou-se a correlação \textbf{dentro} de \emph{m} (12
  obs.: 6 anos × 2 etapas). Trata-se de diagnóstico de co-tendência
  intra-unidade --- não replica a associação transversal nem constitui
  validação externa.
\item
  \textbf{Poder e equivalência.} O mínimo efeito detectável (MDES) segue
  Bloom (1995): com α=5\% bilateral e poder de 80\%, MDES\_β =
  (t\_\{0,975;8\} + t\_\{0,80;8\}) · SE CRVE = (2,306 + 0,889) · 1,881 ≈
  \textbf{6,01 pontos de IDEB por unidade de log do PIB} --- equivalente
  a apenas 0,57 ponto por +10\% de PIB real. O teste de equivalência
  TOST (Lakens; Scheel; Isager, 2018) executou dois testes unicaudais
  contra margens SESOI de ±0,5 ponto de IDEB (ancorada em efeitos
  típicos de programas de alfabetização) e, como transparência, ±0,1;
  equivalência só seria declarada com ambos os p unicaudais \textless{}
  0,05.
\item
  \textbf{Análise macro complementar.} Sete correlações de Pearson de
  longo prazo (1960--2026) entre PIB per capita e educação, PISA,
  internet, Gini, insegurança alimentar, além de ANOVA de renda por
  nível educacional, com IC95\% por métodos assintóticos/boot e
  discussão explícita de cointegração e de amostras curtas.
\item
  \textbf{Robustez por ablação de especificação (ablation).} Protocolo
  repetido para defasagens de 0 a 4 anos --- ablação sistemática da
  escolha temporal ---, verificando estabilidade de sinais e magnitudes
  (Seção 4.7).
\end{enumerate}

\subsubsection{3.5 Critériios de invalidação empírica
(falsificabilidade)}\label{crituxe9riios-de-invalidauxe7uxe3o-empuxedrica-falsificabilidade}

H1 seria falsificada se o IC95\% da correlação between incluísse zero ou
se r \textless{} 0,3. H2 seria falsificada se o valor-p \emph{wild
cluster bootstrap} fosse inferior a 0,05. H3a seria falsificada se a
proporção de metas cumpridas fosse inferior a 95\%. H3b seria
falsificada se o IC95\% da correlação ganho--renda excluísse zero.
Adicionalmente, a disciplina anti-overclaim exige que qualquer resultado
não significativo venha acompanhado de MDES e TOST: \textbf{ausência de
significância sem dimensionamento de poder não autoriza leitura de
nulidade} --- classifica-se como precisão insuficiente.

\subsubsection{3.6 Aprovação ética e
reprodutibilidade}\label{aprovauxe7uxe3o-uxe9tica-e-reprodutibilidade}

O estudo utiliza exclusivamente dados secundários públicos, agregados e
sem identificadores individuais, o que dispensa apreciação por Comitê de
Ética em Pesquisa (registro de isenção documentado com as fontes). A
reprodutibilidade é garantida por: manifesto SHA-256 de todos os
insumos; scripts versionados de coleta e de análise (listados no
Apêndice A); sementes aleatórias fixas (42 para bootstraps; 9.999
réplicas Rademacher pré-especificadas); ambiente congelado
(\texttt{requirements.txt}, \texttt{environment.yml}); e saídas
completas em JSON (\texttt{outputs/expanded/resultados\_r428.json},
\texttt{provenance\_r428.json}). As referências seguem a \textbf{ABNT
NBR 6023}. Qualquer terceiro pode regenerar cada tabela e figura a
partir dos hashes registrados.

\subsection{4. Resultados}\label{resultados}

\subsubsection{4.1 Escala macro: correlações de longo prazo e o platô de
14
anos}\label{escala-macro-correlauxe7uxf5es-de-longo-prazo-e-o-platuxf4-de-14-anos}

Na janela 1960--2026, o PIB per capita brasileiro multiplicou-se por
cerca de 3,6 entre 1960 e 1980, estagnou na década perdida, retomou
trajetória no Plano Real, cresceu 255\% entre 2000 e 2011 e, desde
então, não recuperou o pico histórico: a média de 2012--2024 ficou
26,4\% abaixo de US\$ 13.397. Essa dinâmica em degraus --- expansão,
colapso, recuperação incompleta --- é o retrato quantitativo da
armadilha da renda média no país.

A Tabela 1 consolida as sete correlações de longo prazo que
operacionalizam a hipótese do ecossistema educacional. Seis são
significativas ao nível de 1\% e apresentam magnitudes elevadas: anos de
escolaridade (r=0,86) e, sobretudo, proficiência PISA (r=0,95) dominam a
associação com o PIB per capita; conectividade (internet--PIB entre 0,66
e 0,71 conforme janela) e a relação inversa entre desigualdade e
educação (Gini--educação entre −0,94 e −0,95) completam o quadro de
complementaridades. A única exceção é informativa: PIB e insegurança
alimentar grave não apresentam associação discernível na amostra
(r=−0,39; IC95\% {[}−0,71; 0,09{]}; p=0,155), coerente com a saída
recente do país do Mapa da Fome ter ocorrido de forma relativamente
autônoma do ciclo de renda. A ANOVA complementar --- renda média por
nível educacional atingido --- produz F=310,5 (p\textless0,001) com
η²=0,889; convém reiterar que o cálculo utiliza médias de grupo, não
indivíduos, servindo como descrição de gradiente, não como estimativa
individual.

\textbf{Tabela 1} -- Correlações de Pearson de longo prazo (1960--2026)
e ANOVA educacional --- escala macro.

Par analisado

r

IC95\%

p

Leitura

Anos de escolaridade × PIB pc

,86

{{[}0,71; 0,93{]}}

\textless0,001

Significativa

PISA (proficiência) × PIB pc

,95

{{[}0,88; 0,98{]}}

\textless0,001

Significativa

Uso de internet × PIB pc

,66--0,71

---

\textless0,01

Significativa

Gini × anos de escolaridade

−0,94 a −0,95

---

\textless0,001

Significativa (inversa)

PIB pc × insegurança alimentar grave

−0,39

{{[}−0,71; 0,09{]}}

,155

Não significativa

Renda por nível educacional (ANOVA)

F=310,5

η²=0,889

\textless0,001

Gradiente entre grupos

\emph{Fonte:} elaboração própria com World Bank Open Data, IBGE,
OCDE/PISA, FAO, USPTO e CONAB. IC95\% por normalidade assintótica;
séries curtas discutidas na Seção 6.

Três ressalvas qualificam o bloco macro sem anular sua função
descritiva. Primeiro, séries temporais longas podem compartilhar
tendências estocásticas comuns (cointegração não testada com VECM aqui);
segundo, amostras de 5--15 pontos amplificam sensibilidade a outliers;
terceiro, ausência de controles multivariados impede isolamento de
contributos marginais. Ainda assim, a hierarquia das correlações ---
proficiência acima de escolaridade; desigualdade fortemente associada à
educação --- é consistente com a tese do capital cognitivo e fundamenta
o teste micro das próximas seções: se aprendizagem é o canal relevante,
políticas que elevem IDEB em contextos pobres constituem ruptura
potencial da armadilha, mesmo sem choque de renda.

\subsubsection{4.2 Escala micro: panorama descritivo da
microrregião}\label{escala-micro-panorama-descritivo-da-microrregiuxe3o}

A Tabela 2 apresenta os nove municípios com população censitária, renda
média do responsável (Censo 2022) e os extremos da trajetória do IDEB em
anos iniciais na rede municipal. Duas regularidades saltam à vista.
Primeira, a amplitude de renda é moderada: de R\$ 1.045,85 (Ararendá) a
R\$ 1.676,70 (Crateús) --- razão máxima inferior a 1,7 ---, o que
limita, mas não elimina, o poder de detecção de gradientes transversais.
Segunda, o ganho 2007→2025 é universalmente positivo e de grande
magnitude, variando de +3,2 (Quiterianópolis) a +6,8 pontos (Monsenhor
Tabosa), com média de 5,93 (dp=1,15; mediana 6,2): todos os municípios
mais que dobraram o IDEB inicial.

\textbf{Tabela 2} -- Municípios do Sertão de Crateús: população, renda
censitária e trajetória IDEB (anos iniciais, rede municipal).

\begin{longtable}[]{@{}
  >{\raggedright\arraybackslash}p{(\linewidth - 10\tabcolsep) * \real{0.2250}}
  >{\raggedleft\arraybackslash}p{(\linewidth - 10\tabcolsep) * \real{0.1375}}
  >{\raggedleft\arraybackslash}p{(\linewidth - 10\tabcolsep) * \real{0.2375}}
  >{\raggedleft\arraybackslash}p{(\linewidth - 10\tabcolsep) * \real{0.1375}}
  >{\raggedleft\arraybackslash}p{(\linewidth - 10\tabcolsep) * \real{0.1375}}
  >{\raggedleft\arraybackslash}p{(\linewidth - 10\tabcolsep) * \real{0.1250}}@{}}
\toprule\noalign{}
\begin{minipage}[b]{\linewidth}\raggedright
Município
\end{minipage} & \begin{minipage}[b]{\linewidth}\raggedleft
População
\end{minipage} & \begin{minipage}[b]{\linewidth}\raggedleft
Renda resp. (R\$)
\end{minipage} & \begin{minipage}[b]{\linewidth}\raggedleft
IDEB 2007
\end{minipage} & \begin{minipage}[b]{\linewidth}\raggedleft
IDEB 2025
\end{minipage} & \begin{minipage}[b]{\linewidth}\raggedleft
Ganho
\end{minipage} \\
\midrule\noalign{}
\endhead
\bottomrule\noalign{}
\endlastfoot
Monsenhor Tabosa & 17.101 & 1.183,07 & 2,6 & 9,4 & \textbf{+6,8} \\
Ipaporanga & 11.545 & 1.122,53 & 3,2 & 9,9 & \textbf{+6,7} \\
Nova Russas & 30.583 & 1.381,34 & 3,0 & 9,7 & \textbf{+6,7} \\
Ararendá & 11.069 & 1.045,85 & 3,4 & 9,8 & \textbf{+6,4} \\
Novo Oriente & 27.495 & 1.172,89 & 3,4 & 9,6 & \textbf{+6,2} \\
Tamboril & 24.744 & 1.149,33 & 3,1 & 9,3 & \textbf{+6,2} \\
Independência & 24.024 & 1.223,61 & 3,7 & 9,8 & \textbf{+6,1} \\
Crateús & 76.215 & 1.676,70 & 3,4 & 8,5 & \textbf{+5,1} \\
Quiterianópolis & 20.209 & 1.132,86 & 3,6 & 6,8 & \textbf{+3,2} \\
\end{longtable}

\emph{Fonte:} INEP/IDEB 2025 (rede municipal, anos iniciais); IBGE Censo
Demográfico 2022. População = moradores em domicílios particulares;
renda = rendimento nominal médio mensal das pessoas responsáveis por
domicílios (V06004).

A Figura 1 situa essa trajetória no contexto estadual e nacional:
partindo de patamar inferior à média cearense em 2005 (≈3,1 contra 3,5),
a microrregião cruza a curva estadual por volta de 2015 e, em
2023--2025, consolida-se acima dela --- 9,4 contra 7,1 em 2023, por
exemplo ---, enquanto a média municipal brasileira evolui de forma bem
mais gradual (≈5,9 em 2023). A Figura 2 decompõe o agregado em
trajetórias individuais e revela dois fatos importantes para as seções
inferenciais: (i) a inflexão comum por volta de 2013--2015 ---
coincidente com a maturação do PAIC --- produz forte co-tendência entre
municípios; (ii) a dispersão entre trajetórias diminui visivelmente após
2017, com sete dos nove municípios entre 9,3 e 9,9 em 2025.
Quiterianópolis é o outlier descendente relativo (6,8), e Crateús --- o
município mais rico e populoso --- apresenta o segundo menor nível final
(8,5), pré-figurando a ausência de gradiente de renda que será testada
formalmente.

\begin{figure}[htbp]
\centering
\includegraphics[width=\linewidth]{figuras/fig01_evolucao_ideb.png}
\caption{IDEB anos iniciais, rede municipal: média do Sertão de Crateús (9 municípios), Ceará e Brasil, 2005–2025. Fonte: INEP/IDEB 2025; elaboração própria.}
\end{figure}

\emph{Figura 1 -- IDEB anos iniciais, rede municipal: média do Sertão de
Crateús (9 municípios), Ceará e Brasil, 2005--2025. Fonte: INEP/IDEB
2025; elaboração própria.}

\begin{figure}[htbp]
\centering
\includegraphics[width=\linewidth]{figuras/fig02_spaghetti_municipios.png}
\caption{Trajetórias individuais de IDEB (anos iniciais, rede municipal), 2005–2025, por município. Fonte: INEP/IDEB 2025; elaboração própria.}
\end{figure}

\emph{Figura 2 -- Trajetórias individuais de IDEB (anos iniciais, rede
municipal), 2005--2025, por município. Fonte: INEP/IDEB 2025; elaboração
própria.}

A Tabela 3 consolida as estatísticas descritivas dos painéis de
estimação nas três escalas. Note-se que o IDEB médio da microrregião no
painel (6,16; dp=1,51) supera o cearense (5,53) e o brasileiro (5,02),
enquanto o log do PIB per capita real médio (9,136 ≈ R\$ 9.300) situa-se
abaixo do estadual e muito abaixo do nacional (9,89 ≈ R\$ 19.700) --- a
assimetria educacional-positiva e econômica-negativa que motiva todo o
teste subsequente.

\textbf{Tabela 3} -- Estatísticas descritivas dos painéis de estimação
(IDEB 2013--2023 pareado a ln PIB pc real t−2).

Escala

n (painel)

IDEB média (dp)

IDEB {{[}mín; máx{]}}

ln PIB pc real média (dp)

Sertão de Crateús

,16 (1,51)

,9; 9,6

,136 (0,144)

Ceará

,53 (1,07)

,0; 10,0

,315 (0,401)

Brasil

,02 (1,13)

,5; 10,0

,892 (0,733)

\emph{Fonte:} INEP/IDEB; IBGE PIB dos Municípios deflacionado pelo IPCA
(R\$ de 2021). Missingness de IDEB: 0,5\% (micro), 1,3\% (CE), 26,2\%
(BR).

\subsubsection{4.3 Associação entre municípios (H1): não-detecção
transversal}\label{associauxe7uxe3o-entre-municuxedpios-h1-nuxe3o-detecuxe7uxe3o-transversal}

O estimador between --- correlação entre médias municipais do painel ---
produz \textbf{r=−0,237} (IC95\% bootstrap {[}−0,768; 0,501{]}; p=0,516;
5.000 réplicas, seed 42). A Figura 3 visualiza as nove médias municipais
no plano (log PIB real; IDEB) com reta de mínimos quadrados: a
inclinação levemente negativa é visualmente tênue e dominada pela
posição relativa de Crateús (maior log-renda, IDEB intermediário-baixo)
e do par Monsenhor Tabosa/Nova Russas (rendas intermediárias, IDEB
alto). Três observações qualificam a leitura. Primeira, o sinal negativo
--- contrário ao gradiente nacional --- já sugere que o mecanismo
agregado não opera nesta escala homogênea. Segunda, a largura do IC95\%
(quase 1,3 ponto na escala de correlação) traduz diretamente G=9:
não-detecção aqui significa falta de evidência, não evidência de
ausência --- daí a obrigatoriedade do aparato de MDES/TOST. Terceira, o
jackknife por município mostrou estabilidade qualitativa do coeficiente
(sem reversão de sinal induzida por exclusão simples), descartando
dependência de um único ponto alavancado.

\begin{figure}[htbp]
\centering
\includegraphics[width=\linewidth]{figuras/fig03_scatter_between.png}
\caption{Associação transversal entre municípios: médias do painel (2013–2023) de ln PIB pc real e IDEB; r\_between=−0,24 (bootstrap por cluster, p=0,52). Fonte: INEP; IBGE; elaboração própria.}
\end{figure}

\emph{Figura 3 -- Associação transversal entre municípios: médias do
painel (2013--2023) de ln PIB pc real e IDEB; r\_between=−0,24
(bootstrap por cluster, p=0,52). Fonte: INEP; IBGE; elaboração própria.}

Em síntese, \textbf{H1 não encontra apoio}: a correlação transversal na
microrregião não apenas falha no limiar pré-registrado (r≥0,3) como
aponta para o sentido oposto ao nacional, com incerteza compatível tanto
com efeitos positivos modestos quanto com negativos expressivos. A
implicação metodológica é imediata --- qualquer afirmação de gradiente
renda-IDEB construída a partir do agregado nacional não se transfere
automaticamente para esta microrregião.

\subsubsection{4.4 Associação dentro dos municípios (H2): quatro
evidências
convergentes}\label{associauxe7uxe3o-dentro-dos-municuxedpios-h2-quatro-eviduxeancias-convergentes}

A Tabela 4 resume o quarteto de estimadores dedicados à dimensão
temporal. Nas \textbf{correlações pooled} (todas as 108 observações),
r=0,159 (IC95\% bootstrap {[}−0,058; 0,424{]}; p=0,141) ---
numericamente positivo, porém indistinguível de zero e muito distante do
pooled nacional (0,415). Nas \textbf{primeiras diferenças} (n=90 pares
consecutivos intra município-etapa), r=−0,071 (p=0,504): ano em que a
renda local acelera não é ano em que o IDEB acelera. No \textbf{modelo
de efeitos fixos} de município×etapa e ano --- que compara cada
município-etapa consigo mesmo, líquidos choques comuns de calendário ---
o coeficiente de ln PIB real t−2 é β=2,607: em elasticidade, +10\% de
PIB per capita real associar-se-ia a +0,25 ponto de IDEB (2,607 × ln
1,1), mas com enorme imprecisão: SE CRVE=1,881, IC95\% {[}−1,73;
6,95{]}, t=1,39 com 8 graus de liberdade, p=0,203. O procedimento
decisivo, \textbf{wild cluster bootstrap Rademacher sob H0} (9.999
réplicas), confirma a não-rejeição: p=0,230. O erro-padrão homocedástico
(1,284; p=0,046) é reportado apenas como transparência --- e ilustra
exatamente o risco que a correção de pequena amostra existe para evitar:
ignorar a estrutura de clusters produziria um falso positivo de 5\%.

\textbf{Tabela 4} -- Inferência micro sobre a associação renda--IDEB
(painel 2013--2023, n=108; G=9 clusters).

Estimador

Estatística

IC95\% / SE

p

Decisão sobre H2

Between (médias municipais)

r=−0,237

{{[}−0,768; 0,501{]}}

,516

Sem associação detectável

Pooled (Pearson)

r=0,159

{{[}−0,058; 0,424{]}}

,141

Sem associação detectável

Primeiras diferenças

r=−0,071

--- (nominal)

,504

Sem co-movimento

Efeitos fixos (mun.×etapa + ano)

β=2,607

EP cluster 1,881 {{[}−1,73; 6,95{]}}

,203 (wild 0,230)

Indistinguível de zero

\emph{Fonte:} elaboração própria. Bootstraps com 5 mil réplicas entre e
pooled, e 9.999 réplicas wild; semente fixa 42; teste t com G−1 graus de
liberdade.

A Figura 6 (Seção 4.7) complementa com o perfil de defasagens 0--4:
nenhuma janela produz correlação pooled superior a 0,26 ou primeira
diferença que se afaste consistentemente de zero --- a conclusão não é
artefato da escolha de lag.

\textbf{Veredito sobre H2:} os quatro estimadores convergem para a mesma
leitura --- \textbf{não há associação temporal discernível entre renda
municipal e IDEB} na microrregião. Contudo, a honestidade inferencial
obriga ao passo seguinte.

\subsubsection{4.5 Poder, equivalência e validação
interna}\label{poder-equivaluxeancia-e-validauxe7uxe3o-interna}

Sem dimensionamento, ``não significativo'' seria retoricamente ambíguo.
O MDES cluster-robusto é de \textbf{6,01 pontos de IDEB por unidade de
log do PIB} --- o estudo só conseguiria detectar, com 80\% de poder,
efeitos equivalentes a quase triplicar a renda local; convertido para a
métrica de política pública, trata-se de \textbf{0,57 ponto de IDEB por
+10\% de PIB per capita real} (versão homocedástica: 3,63/log-unit). Em
outras palavras: efeitos plausíveis de curto prazo da economia municipal
sobre aprendizagem estão abaixo do piso de detecção deste desenho.
Consequentemente, o \textbf{TOST com margem substantiva de ±0,5 ponto}
--- ordem de grandeza dos impactos típicos de programas de alfabetização
--- não estabelece equivalência (p=0,852; margem auxiliar ±0,1:
p=0,890). A classificação correta do achado é, portanto:
\textbf{precisão insuficiente para sustentar efeito ou equivalência} ---
formulação que protege o artigo tanto do overclaim causal quanto da
pseudo-confirmação de irrelevância da renda.

A validação interna por LOOCV adiciona nuance valiosa. Quando cada
município é retirado e a correlação de níveis \textbf{dentro dele} é
computada, os valores são uniformemente positivos --- média 0,443;
mediana 0,459; dp=0,203; mínimo 0,08 (Quiterianópolis); máximo 0,72
(Ararendá); 100\% positivos (Tabela 5). À primeira vista paradoxal
frente aos betas nulos, o resultado é metodologicamente coerente: dentro
de cada município, IDEB e PIB \textbf{co-tendenciam} ao longo de quinze
anos (ambos crescem com o tempo), mas isso não implica associação
condicional de curto prazo --- exatamente o que os efeitos fixos com
dummies de ano isolam. A co-tendência positiva intra-unidade convive com
ausência de gradiente entre unidades e com nulidade condicional: o trio
entre/dentro/co-tendência descreve um caso em que o avanço educacional
compartilhado domina a variância temporal comum.

\textbf{Tabela 5} -- LOOCV: correlação de níveis dentro de cada
município retido (12 obs. cada).

\begin{longtable}[]{@{}llll@{}}
\toprule\noalign{}
Município & r dentro & Município & r dentro \\
\midrule\noalign{}
\endhead
\bottomrule\noalign{}
\endlastfoot
Quiterianópolis & 0,08 & Ipaporanga & 0,46 \\
Novo Oriente & 0,27 & Crateús & 0,51 \\
Independência & 0,30 & Nova Russas & 0,65 \\
Monsenhor Tabosa & 0,32 & Tamboril & 0,67 \\
Ararendá & 0,72 & Média (dp) & 0,44 (0,20) \\
\end{longtable}

\emph{Fonte:} elaboração própria (validação interna; ver nota conceitual
em 3.4, item 5).

\subsubsection{4.6 Metas do INEP (H3a) e ganho educacional
(H3b)}\label{metas-do-inep-h3a-e-ganho-educacional-h3b}

\textbf{Cumprimento de metas.} Dos 108 pares município×etapa×ano
avaliáveis entre 2011 e 2021, \textbf{107 (99,1\%) atingiram ou
superaram a meta projetada pelo INEP no mesmo ano de referência}. A
única exceção ocorreu em 2013, anos iniciais de Ipaporanga: IDEB
observado de 4,7 contra meta de 5,0 --- déficit de 0,3 ponto em uma
única célula, seguido de recuperação imediata (6,8 em 2015, acima da
meta daquele ano). A Figura 4 mostra os percentuais anuais: 100\% em
2011, 94,4\% em 2013 e 100\% em todos os anos subsequentes. Sob a ótica
da accountability, o desempenho é notável: em região de baixa renda, o
sistema municipal manteve compromissos públicos de aprendizagem por
cinco ciclos consecutivos com taxa de sucesso próxima da totalidade ---
padrão difícil de conciliar com a hipótese de que a renda local é
pré-condição de desempenho.

\begin{figure}[htbp]
\centering
\includegraphics[width=\linewidth]{figuras/fig04_metas_por_ano.png}
\caption{Percentual de metas do INEP cumpridas no mesmo ano de referência (municipal, AI e AF; n=108 pares 2007–2021). Fonte: INEP/IDEB; elaboração própria.}
\end{figure}

\emph{Figura 4 -- Percentual de metas do INEP cumpridas no mesmo ano de
referência (municipal, AI e AF; n=108 pares 2007--2021). Fonte:
INEP/IDEB; elaboração própria.}

\textbf{Ganho e sua (não) associação com renda.} O ganho médio de IDEB
em anos iniciais entre 2007 e 2025 foi de \textbf{+5,93 pontos}
(dp=1,15; mediana 6,2; amplitude 3,2--6,8), com oito dos nove municípios
ganhando entre +5,1 e +6,8 e Quiterianópolis como caso-atípico (+3,2,
ainda assim representando quase duplicação do ponto de partida). A
correlação entre esse ganho e a renda censitária municipal é
\textbf{r=−0,136 (p=0,726; n=9)}: se algo, a direção é levemente
invertida --- os municípios marginalmente mais ricos não ganharam mais.
A Figura 5 exibe o gráfico correspondente com ajuste linear quase plano.
Combinados, H3a e H3b descrevem um sistema que (i) cumpre virtualmente
todas as metas externas e (ii) cuja velocidade de avanço é ortogonal à
renda local --- assinatura empírica de mecanismo
institucional-pedagógico, e não econômico.

\begin{figure}[htbp]
\centering
\includegraphics[width=\linewidth]{figuras/fig05_ganho_vs_renda.png}
\caption{Ganho de IDEB (AI, 2007→2025) versus renda média do responsável (Censo 2022); r=−0,14 (p=0,73; n=9). Fonte: INEP; IBGE Censo 2022; elaboração própria.}
\end{figure}

\emph{Figura 5 -- Ganho de IDEB (AI, 2007→2025) versus renda média do
responsável (Censo 2022); r=−0,14 (p=0,73; n=9). Fonte: INEP; IBGE Censo
2022; elaboração própria.}

\subsubsection{4.7 Robustez: defasagens alternativas, poder e
especificações}\label{robustez-defasagens-alternativas-poder-e-especificauxe7uxf5es}

O perfil de defasagens (Tabela 6; Figura 6) mostra que nenhuma escolha
razoável de lag altera as conclusões: as correlações pooled flutuam
entre 0,159 e 0,256 sem padrão monotônico, e as primeiras diferenças
permanecem em banda estreita ao redor de zero (−0,07 a 0,18, todas
nominais não significativas). A defasagem principal de 2 anos não é um
ponto favorável isolado, mas parte de uma superfície plana ---
exatamente o que se espera quando não há mecanismo temporal robusto
conectando as séries.

\textbf{Tabela 6} -- Perfil de robustez por defasagem do PIB real
(painel micro, n=108 por lag).

\begin{longtable}[]{@{}lll@{}}
\toprule\noalign{}
Lag (anos) & r níveis (pooled) & r primeiras diferenças \\
\midrule\noalign{}
\endhead
\bottomrule\noalign{}
\endlastfoot
& 0,201 & 0,053 \\
1 & 0,256 & 0,171 \\
\textbf{2 (principal)} & \textbf{0,159} & \textbf{−0,071} \\
3 & 0,222 & −0,033 \\
4 & 0,197 & 0,176 \\
\end{longtable}

\emph{Fonte:} elaboração própria; mesmas sementes e réplicas do
protocolo principal.

\begin{figure}[htbp]
\centering
\includegraphics[width=\linewidth]{figuras/fig06_perfil_lags.png}
\caption{Perfil de correlações por defasagem do PIB real (0–4 anos), níveis pooled e primeiras diferenças. Fonte: elaboração própria.}
\end{figure}

\emph{Figura 6 -- Perfil de correlações por defasagem do PIB real (0--4
anos), níveis pooled e primeiras diferenças. Fonte: elaboração própria.}

\subsubsection{4.8 Triangulação de escalas: Crateús, Ceará e
Brasil}\label{triangulauxe7uxe3o-de-escalas-crateuxfas-cearuxe1-e-brasil}

A Tabela 7 alinha as três escalas sob estimadores idênticos. O gradiente
transversal cresce com a escala --- de −0,24 (Crateús) a 0,03 (Ceará) e
0,49 (Brasil) ---, enquanto a dimensão temporal vai de indistinguível de
zero (β=2,61, p wild=0,23; Ceará β≈0,00, p=0,997) a pequena porém
precisa no país (β=0,114; SE CRVE=0,024; p\textless0,001). Essa
assimetria é analiticamente rica: (i) confirma que o gradiente nacional
conhecido é composto por comparações entre unidades profundamente
heterogêneas; (ii) mostra que mesmo no agregado estadual cearense ---
onde 184 redes convivem sob o PAIC --- o gradiente desaparece, sugerindo
que a coordenação comprime diferenças entre municípios; e (iii) fornece,
na célula nacional, o \textbf{baseline} de comparação com o estado da
arte --- onde a coordenação é ausente ou heterogênea, a componente
temporal da renda reaparece com sinal positivo e preciso (β=0,114;
p\textless0,001), ainda que modesta por ano de defasagem.

\textbf{Tabela 7} -- Triangulação de escalas: between, efeitos fixos e
dimensões de amostra.

\begin{table}[htbp]
\centering
\footnotesize
\setlength{\tabcolsep}{3pt}
\begin{tabular}{@{}lrrlll@{}}
\hline
Escala & G & n & Between r [IC] & FE $\beta$ (EP) & p / wild \\
\hline
Sertão de Crateús & 9 & 108 & $-0,237$ [$-0,768$; 0,501] & 2,607 (1,881) & 0,203 / 0,230 \\
Ceará & 184 & 2.202 & 0,027 [$-0,103$; 0,155] & $-0,0006$ (0,138) & 0,997 / --- \\
Brasil & 5.360 & 48.695 & 0,489 [0,471; 0,508] & 0,114 (0,024) & $<$0,001 / --- \\
\hline
\end{tabular}
\end{table}

\emph{Fonte:} INEP/IDEB; IBGE; protocolo idêntico ao da Seção 3. Wild
bootstrap executado apenas na micro (G\textless40); nas escalas maiores,
t(G−1) é adequado.

\subsubsection{4.9 Síntese dos achados}\label{suxedntese-dos-achados}

\begin{enumerate}
\def\labelenumi{(\roman{enumi})}
\tightlist
\item
  Na macroescala, educação --- especialmente proficiência --- é o
  correlato mais forte do PIB per capita, inserida em ecossistema que
  inclui conectividade, IA e desigualdade. (ii) No Sertão de Crateús,
  não há associação detectável entre renda municipal e IDEB, nem entre
  municípios, nem dentro deles ao longo de seis ondas, sob protocolo
  conservador com wild bootstrap. (iii) O poder do desenho é
  explicitamente insuficiente para efeitos plausíveis (MDES=0,57 ponto
  por +10\% de PIB), e a equivalência com efeito nulo não pode ser
  afirmada (TOST p=0,85). (iv) A co-tendência interna é universalmente
  positiva (LOOCV 0,08--0,72), refletindo avanço conjunto no tempo. (v)
  Metas do INEP foram cumpridas em 99,1\% das células, com ganho médio
  de +5,93 pontos desacoplado da renda (r=−0,14). (vi) A triangulação
  indica que o gradiente renda--IDEB é fenômeno de escala nacional,
  comprimido pela coordenação estadual e ausente na microrregião.
\end{enumerate}

\subsection{5. Discussão}\label{discussuxe3o}

\subsubsection{5.1 Interpretação teórica: capital cognitivo antes da
renda}\label{interpretauxe7uxe3o-teuxf3rica-capital-cognitivo-antes-da-renda}

Os resultados articulam-se com precisão ao arcabouço de Hanushek e
Woessmann (2015): aquilo que liga educação a desenvolvimento é
aprendizagem efetiva, e a evidência cearense mostra que aprendizagem
básica pode avançar substancialmente \textbf{antes} --- e
independentemente --- da renda municipal. No plano da armadilha da renda
média, isso reordena a sequência implícita de políticas: em vez de
tratar melhoria escolar como luxo derivado de crescimento, o caso
Crateús demonstra a viabilidade inversa, na qual a fase de capital
humano básico é conduzida pela engenharia institucional da política
educacional. A ausência de gradiente transversal local, somada ao ganho
médio de quase seis pontos, equivale a dizer que o fator que explica o
IDEB não reside na dotação econômica municipal, mas no pacote
pedagógico-governativo compartilhado.

\subsubsection{5.2 Mecanismos institucionais: o que o PAIC parece
fazer}\label{mecanismos-institucionais-o-que-o-paic-parece-fazer}

Embora este desenho não identifique canais causais, a cronologia é
sugestiva: a inflexão coletiva das trajetórias (Figura 2) coincide com a
maturação do PAIC --- formação continuada de professores
alfabetizadores, materiais estruturados, avaliação externa anual de
fluência no 2º ano e governança de resultados. Três características
desse arranjo são teoricamente relevantes: \textbf{regularidade} (ciclos
curtos e previsíveis de meta-avaliação), \textbf{complementaridade
federativa} (estado fornece bens públicos pedagógicos que municípios
pequenos não teriam escala para produzir) e \textbf{visibilidade}
(rankings e metas tornam desempenho informação pública acionável). A
literatura de accountability alerta, com razão, para custos potenciais
--- ensino para o teste, estreitamento curricular, gaming de fluxo
(Afonso, 2012) ---; os dados aqui apresentados não permitem medir esses
custos, apenas registrar que o indicador-síntese avançou de forma
generalizada e sustentada por mais de uma década, com metas cumpridas em
99,1\% das células.

\subsubsection{5.3 Implicações para a armadilha da renda média e para a
política
pública}\label{implicauxe7uxf5es-para-a-armadilha-da-renda-muxe9dia-e-para-a-poluxedtica-puxfablica}

A implicação central é de sequenciamento: \textbf{municípios pobres não
precisam esperar o ciclo econômico para produzir a etapa de capital
humano básico}. Num país com 14 anos de platô macroeconômico, essa é uma
boa notícia operacional --- a janela educacional não está fechada pela
renda local. Simultaneamente, os resultados macro temperam o otimismo:
IDEB elevado é condição habilitante, não gatilho suficiente de ruptura
de renda; conectividade, inovação industrial e distribuição de renda
permanecem complementares. Para o gestor municipal, a mensagem prática é
dupla: (i) metas de IDEB são alcançáveis no horizonte de dois a três
ciclos de política (4--6 anos) mesmo com orçamento restrito, como
demonstram oito municípios que partiram de 2,6--3,7 e chegaram a
8,5--9,9; (ii) a variável de atenção não deve ser a conjuntura
econômica, mas a constância do pacote pedagógico --- exatamente a
dimensão sob controle administrativo. Para o formulador
estadual/nacional, o caso reforça o valor dos regimes de colaboração com
avaliação externa regular como tecnologia de difusão de aprendizagem em
redes pequenas.

\subsubsection{5.4 Posição no debate
internacional}\label{posiuxe7uxe3o-no-debate-internacional}

Comparado à evidência multinível clássica --- em que o nível
socioeconômico explica parcela relevante da variância entre escolas e
alunos ---, o achado municipal cearense não o contradiz: opera noutra
fronteira de variação. Entre nove territórios institucionalmente
homogêneos e socioeconomicamente compactos, a variância explicável por
renda colapsa; entre países ou entre redes sem política comum, ela
aflora (Brasil between 0,49). Essa leitura reconcilia os dois blocos de
literatura e sugere que a questão empiricamente fértil não é ``quanto a
renda explica?'', mas ``\textbf{sob quais arranjos institucionais a
renda deixa de explicar?}'' --- reformulação à qual este artigo offering
um caso de referência e um protocolo replicável.

\subsection{6. Limitações}\label{limitauxe7uxf5es}

\textbf{Precisão e poder.} Com G=9, o IC95\% da correlação between é
amplo ({[}−0,77; 0,50{]}) e o MDES é elevado (6,01/log-unit; 0,57 ponto
por +10\% de PIB): efeitos verdadeiros moderados poderiam passar
despercebidos. O TOST não estabeleceu equivalência; portanto, o estudo
\textbf{não} demonstra que a renda é irrelevante --- demonstra que seu
efeito local, se existir, é menor do que o detectável neste desenho.

\textbf{Identificação.} O desenho é observacional; não há instrumento
nem quasi-experimento. Coeficientes capturam associação condicional;
choques de renda correlacionados a choques pedagógicos não podem ser
descartados conceitualmente, ainda que a ortogonalidade descritiva (PD≈0
em cinco lags) torne tais confundidores pouco prováveis de dominar.

\textbf{Validade de constructo do IDEB.} Como todo indicador síntese, o
IDEB é sensível à composição fluxo×proficiência e sujeito às críticas de
accountability (estreitamento curricular; gestão de matrícula). Avanços
do IDEB nem sempre equivalem ponto a ponto a avanços de aprendizagem
mais amplos; a decomposição municipal por componente não estava
disponível de forma auditável para toda a série.

\textbf{Medida de renda.} PIB per capita municipal é medida de atividade
produtiva, não de bem-estar domiciliar; a renda censitária do
responsável é corte único (2022). O IPCE comercial, originalmente
desejado, carece de auditabilidade pública. Diferentes proxies de nível
socioeconômico poderiam produzir magnitudes distintas, ainda que
improvavelmente revertam o quadro de não-detecção dado o MDES.

\textbf{Macro: séries curtas e cointegração.} As correlações de
1960--2026 usam 5--15 pontos efetivos por par e não tratam raízes
unitárias/cointegração (VECM); devem ser lidas como descrições de longo
prazo, não como elasticidades estruturais. A ANOVA educacional usa
médias de grupo (η² não é efeito individual).

\textbf{Generalização.} O Sertão de Crateús é um caso-limite de
coordenação estadual eficaz. A external validity para municípios fora de
regimes de colaboração maduros é limitada por construção --- e é
justamente essa a proposição substantiva do artigo.

\subsection{7. Conclusão}\label{conclusuxe3o}

Este artigo perguntou se a educação --- a variável mais fortemente
correlacionada ao desenvolvimento na escala macro --- pode ser produzida
na escala municipal independentemente da renda local, e respondeu
afirmativamente com evidência de um caso crítico do Semiárido
brasileiro. Na macroescala, proficiência (PISA--PIB r=0,95) e
escolaridade (r=0,86) dominam o quadro correlacional de longo prazo,
inseridas em um ecossistema que inclui conectividade digital, adoção
tecnológica e desigualdade; a estagnação brasileira de 14 anos abaixo do
pico de 2011 mostra que nenhum desses fatores, sozinho, rompeu a
armadilha. Na microescala do Sertão de Crateús, nove municípios de baixa
renda elevaram o IDEB de anos iniciais em média 5,93 pontos entre 2007 e
2025, cumpriram 107 de 108 metas do INEP no ano de referência, e --- sob
bootstrap por cluster, CRVE corrigido, wild bootstrap, MDES/TOST e LOOCV
--- não apresentaram associação discernível entre renda municipal e
desempenho, nem entre municípios (r=−0,24), nem dentro deles (β=2,61; p
wild=0,23).

Esta contribuição é \textbf{inédita} e \textbf{original} em três planos
--- e a \textbf{metodologia} proposta é replicável: (i) empírica, ao
documentar com protocolo conservador e proveniência criptográfica um
caso de desacoplamento educação--renda na escala municipal; (ii)
metodológica, ao transferir para a pesquisa educacional municipal o
arsenal completo de inferência com poucos clusters (bootstrap de
clusters, CRVE de pequena amostra, wild bootstrap, MDES, TOST, LOOCV),
com critérios de falsificabilidade pré-registrados; e (iii) de política
pública, ao reordenar o sequenciamento da ruptura da armadilha --- a
etapa de capital humano básico pode, e no Ceará foi, conduzida antes e
independentemente do crescimento econômico local, por regimes de
colaboração com avaliação regular.

Pesquisas futuras devem: incorporar decomposições de proficiência e
fluxo do IDEB para rastrear mecanismos; estender o painel a
microrregiões de outros estados com e sem arranjos análogos ao PAIC,
explorando diferenças-em-diferenças entre regimes de colaboração; tratar
explicitamente a cointegração das séries macro (VECM/ARDL) com painéis
de país ampliados; e avaliar desfechos de longo prazo (trajetórias
escolares dos cohorts alfabetizados sob o PAIC) para conectar a ruptura
educacional local a desfechos de produtividade --- o elo que a teoria da
armadilha exige e que este artigo deixa, deliberadamente, como próxima
fronteira.

\subsection{Referências}\label{referuxeancias}

\begingroup\small

AFONSO, A. J. Para uma concetualização alternativa de accountability em
educação. \emph{Educação \& Sociedade}, Campinas, v. 33, n.~119,
p.~471-484, 2012.

AIYAR, S. et al.~\emph{Growth Slowdowns and the Middle-Income Trap}.
Washington, DC: IMF Working Paper WP/13/71, 2013.

ANDREWS, C.; VRIES, W. Pobreza e municipalização no Brasil: expectativas
e desafios. \emph{Cadernos de Pesquisa}, São Paulo, v. 42, n.~147,
p.~1-20, 2012.

BANCO MUNDIAL. \emph{World Development Indicators}: PIB per capita
Brasil 1960--2024. Washington, DC: World Bank Open Data, 2024.
Disponível em: https://data.worldbank.org. Acesso em: 15 ago. 2026.

BANCO MUNDIAL. \emph{Gini index (Brazil)}. Washington, DC: World Bank
Open Data, 2024.

BLOOM, H. S. Minimum detectable effects: a primer. \emph{Journal of
Research on Educational Effectiveness}, v. 1, n.~1, p.~1-10, 1995.

BRASIL. \emph{Constituição da República Federativa do Brasil de 1988}.
Brasília: Senado Federal, 1988.

BRUNELLO, G.; KISS, I. Does high-stakes testing improve educational
outcomes? \emph{Economics of Education Review}, v. 89, p.~102-110, 2022.

CAMERON, A. C.; GELBACH, J. B.; MILLER, D. L. Bootstrap-based
improvements for inference with clustered errors. \emph{The Review of
Economics and Statistics}, v. 90, n.~3, p.~414-427, 2008.

CARNOY, M. et al.~How education policy can affect student learning: an
intranational analysis. \emph{International Journal of Educational
Development}, v. 56, p.~1-10, 2017.

CHERIF, R.; HASANOV, F. \emph{The Leap of the Tiger}: how Malaysia can
escape the middle-income trap. Washington, DC: IMF, 2015.

CHERIF, R.; HASANOV, F. \emph{Pitfalls to avoid when escaping the
middle-income trap}. Washington, DC: IMF, 2024.

CHERIF, R.; HASANOV, F. Escaping the middle-income trap: a conceptual
framework. \emph{Journal of Development Economics}, v. 150, p.~1-15,
2024.

CITO, C.; MARÔCO, J. Desigualdades socioeconômicas no PISA: o caso
brasileiro. \emph{Revista Portuguesa de Educação}, v. 39, n.~1, p.~1-20,
2026.

CONAB. \emph{Perspectivas para a agropecuária}: produtividade agrícola
2024. Brasília: CONAB, 2024.

COSTA, L.; CARNOY, M. The effectiveness of an early-grade literacy
intervention on the cognitive achievement of Brazilian students.
\emph{Educational Evaluation and Policy Analysis}, v. 37, n.~4,
p.~475-489, 2015.

CRUZ, P.; RIBEIRO, C.; BATISTA, N. C. Aprendizagem na idade certa: a
experiência do Ceará. \emph{Revista Brasileira de Estudos Pedagógicos},
Brasília, v. 103, n.~265, p.~1-20, 2022.

DUARTE, N. S. O impacto da pobreza no IDEB: uma análise multinível.
\emph{Ensaio: Avaliação e Políticas Públicas em Educação}, Rio de
Janeiro, v. 21, n.~81, p.~719-734, 2013.

FAO. \emph{The State of Food Security and Nutrition in the World 2024}.
Roma: FAO, 2024.

GILL, I.; KHARAS, H. \emph{An East Asian Renaissance}: ideas for
economic growth. Washington, DC: World Bank, 2007.

GILL, I.; KHARAS, H. \emph{The Middle-Income Trap Turns Ten}.
Washington, DC: World Bank Policy Research Working Paper 7403, 2015.

HANUSHEK, E. A.; WOESSMANN, L. \emph{The Knowledge Capital of Nations}:
education and the economics of growth. Cambridge: MIT Press, 2015.

HANUSHEK, E. A.; WOESSMANN, L. The role of education and skills in
economic development. \emph{Journal of Economic Literature}, v. 60,
n.~2, p.~1-30, 2022.

IBGE. \emph{PIB dos Municípios} 2010--2021. Rio de Janeiro: IBGE, 2023.
Disponível em: ftp://ftp.ibge.gov.br. Acesso em: 15 ago. 2026.

IBGE. \emph{Censo Demográfico 2022}: rendimento nominal mensal das
pessoas responsáveis por domicílios (V06004/V06006). Rio de Janeiro:
IBGE, 2023.

IBGE. \emph{Malha municipal 2023}. Rio de Janeiro: IBGE, 2023.

IBGE. \emph{PINTEC 2022}: Pesquisa de Inovação Tecnológica nas Empresas.
Rio de Janeiro: IBGE, 2024.

IBGE. \emph{PNAD Contínua}: índice de Gini 2024. Rio de Janeiro: IBGE,
2024.

IBGE. \emph{SIDRA}: Censo Demográfico 2022, agregados por municípios.
Rio de Janeiro: IBGE, 2024.

INEP. \emph{IDEB 2005--2025}: resultados e metas, rede municipal.
Brasília: INEP, 2025. Disponível em: https://download.inep.gov.br.
Acesso em: 15 ago. 2026.

INEP. \emph{Nota técnica}: IDEB 2023 --- cálculo e metas. Brasília:
INEP, 2023.

IPEA. \emph{IPCA --- número-índice}. Brasília: IPEADATA, 2023.

IPEA. Séries históricas de PIB e renda municipal. Brasília: IPEA, 2024.

LACRUZ, A. J.; AMÉRICO, B. L.; CARNIEL, F. Indicadores de desempenho
educacional e accountability. \emph{Revista de Administração Pública},
Rio de Janeiro, v. 53, n.~6, p.~1-20, 2019.

LAKENS, D.; SCHEEL, A. M.; ISAGER, P. M. Equivalence testing for
psychological research: a tutorial. \emph{Advances in Methods and
Practices in Psychological Science}, v. 1, n.~2, p.~259-269, 2018.

MACKINNON, J. G.; WEBB, M. D. Wild bootstrap inference for wildly
different cluster sizes. \emph{Journal of Applied Econometrics}, v. 32,
n.~2, p.~233-254, 2017.

MADDISON PROJECT DATABASE. \emph{Historical GDP per capita}, versão
2023. Groningen: University of Groningen, 2023.

OCDE. \emph{PISA 2022 Results}, v. I-II. Paris: OECD Publishing, 2023.

OCDE. \emph{Education at a Glance 2023}: OECD Indicators. Paris: OECD
Publishing, 2023.

PENN WORLD TABLE. \emph{PWT 10.01}. Groningen: University of Groningen,
2023.

SCHNEIDER, M. P.; NARDI, E. L. O IDEB e a construção de um modelo de
accountability na educação básica brasileira. \emph{Revista Brasileira
de Política e Administração da Educação}, v. 30, n.~1, p.~1-20, 2014.

SEGATTO, C.; OLIVEIRA, R.; SILVA, F. Coordenação federativa e resultados
educacionais: o caso cearense. \emph{Revista de Administração Pública},
Rio de Janeiro, v. 58, n.~2, p.~1-20, 2024.

SOARES, J. F.; XAVIER, F. P. Pressupostos educacionais e estatísticos do
IDEB. \emph{Educação \& Sociedade}, Campinas, v. 34, n.~124, p.~903-923,
2013.

UNESCO. \emph{Global Education Monitoring Report 2023}. Paris: UNESCO,
2023.

USPTO. \emph{Patent Statistics Report 2024}. Alexandria: USPTO, 2024.

WORLD BANK. \emph{World Development Indicators} (série completa
1960--2026). Washington, DC: World Bank Open Data, 2026.

\subsection{Apêndice A --- Proveniência, reprodutibilidade e material
suplementar}\label{apuxeandice-a-proveniuxeancia-reprodutibilidade-e-material-suplementar}

\textbf{Manifesto de proveniência.} Cada arquivo-fonte baixado tem URL,
data-hora e hash SHA-256 registrados no manifesto de proveniência do
projeto Crateús-IDEB (pasta data/raw). O deflator IPCA médio anual, com
fator por ano para R\$ de 2021, está na pasta data/processed.

\textbf{Saídas analíticas.} Todos os estatísticos citados no artigo
estão consolidados no arquivo de resultados do ciclo R428 do projeto
Crateús-IDEB (between, pooled, primeiras diferenças, efeitos fixos com
CRVE e wild bootstrap, LOOCV, MDES/TOST, metas por ano, ganhos
municipais e perfil de lags), acompanhados do registro completo de
provenância com sementes e número de réplicas.

\textbf{Scripts.} \texttt{baixar\_dados.py} (coleta com verificação de
hash); \texttt{analise\_crateus.py} e \texttt{analise\_crateus\_r428.py}
(estimações); \texttt{scripts/gerar\_figuras\_tabelas.py} (figuras e
tabelas deste artigo, diretamente dos JSONs processados).

\textbf{Relatório de validação independente.} A validação por banca
simulada multi-periódico (protocolo interno R439:
CAPES/Nature/IEEE/Lancet, com limpeza de gaps) consta de documento
separado, não integrando o corpo do artigo, em conformidade com o padrão
editorial de periódicos.

\textbf{Reprodutibilidade rápida.}

\endgroup

\textbf{Agradecimentos:} à Secretaria de Educação do Estado do Ceará e
às nove secretarias municipais do Sertão de Crateús, pela transparência
dos dados públicos que viabilizaram esta análise.

\end{document}
